\hypertarget{chapter-05-page-113}{}
\subsection{Calderón--Zygmund 奇异积分}
\label{sec:ch5-calderon-zygmund-singular-integrals}

上一节的例子似乎表明, 对任意 Calderón--Zygmund 算子, 至少当 $f\in\mathcal S$ 时应有
\[
 Tf(x)=\lim_{\varepsilon\to0}\int_{|x-y|>\varepsilon}K(x,y)f(y)\,dy.
\]
但这未必成立. 由 \eqref{eq:ch5-standard-kernel-size}, 截断算子 $T_\varepsilon f(x)=\int_{|x-y|>\varepsilon}K(x,y)f(y)\,dy$ 对 $f\in\mathcal S$ 有定义, 可是当 $\varepsilon\to0$ 时极限可能不存在, 或虽存在却不等于 $Tf(x)$. 例~\ref{ex:ch5-complex-power-kernel} 的卷积算子给出第一种情形.

\begin{Proposition}
\label{prop:ch5-truncation-existence-criterion}
极限
\begin{equation}
\label{eq:ch5-generalized-truncation-limit}
 \lim_{\varepsilon\to0}T_\varepsilon f(x)
\end{equation}
对每个 $f\in C_c^\infty$ 几乎处处存在, 当且仅当
\[
 \lim_{\varepsilon\to0}\int_{\varepsilon<|x-y|<1}K(x,y)\,dy
\]
对几乎处处的 $x$ 存在.
\end{Proposition}

极限 \eqref{eq:ch5-generalized-truncation-limit} 存在并不蕴含它等于 $Tf(x)$. 恒等算子 $I$ 是与零核相联系的 Calderón--Zygmund 算子, 但它的截断极限为零. 这个例子也说明算子不由其核唯一决定. 事实上, 任意逐点乘法算子 $Tf(x)=a(x)f(x)$, $a\in L^\infty$, 都与零核相联系. 反之有下述结论, 其证明留给读者.

\begin{Proposition}
\label{prop:ch5-same-kernel-difference}
若两个 Calderón--Zygmund 算子与同一个核相联系, 则它们之差是逐点乘法算子.
\end{Proposition}

必须假设 Calderón--Zygmund 算子在 $L^2$ 上有界. 否则命题~\ref{prop:ch5-same-kernel-difference} 不成立, 例如微分算子的核为零, 但它不是逐点乘法算子.

\hypertarget{chapter-05-page-114}{}
称满足
\begin{equation}
\label{eq:ch5-cz-singular-integral-definition}
 Tf(x)=\lim_{\varepsilon\to0}T_\varepsilon f(x)
\end{equation}
的 Calderón--Zygmund 算子为 Calderón--Zygmund 奇异积分. 为判断 \eqref{eq:ch5-cz-singular-integral-definition} 对 $f\in L^p$ 何时成立, 像处理 Hilbert 变换时那样考察极大算子
\[
 T^*f(x)=\sup_{\varepsilon>0}|T_\varepsilon f(x)|.
\]
由定理~\ref{thm:ch2-closed-ae-convergence-set}, 若 $T^*$ 是弱型 $(p,p)$ 的, 则 $L^p$ 中截断极限几乎处处存在的函数集合是闭集. 所以只要 \eqref{eq:ch5-cz-singular-integral-definition} 对一个稠密子集成立, 它就对每个 $f\in L^p$ 成立.

\begin{Theorem}
\label{thm:ch5-maximal-cz-bounds}
若 $T$ 是 Calderón--Zygmund 算子, 则 $T^*$ 是强型 $(p,p)$ 的, $1<p<\infty$, 并且是弱型 $(1,1)$ 的.
\end{Theorem}

证明依赖于下述类似 Cotlar 不等式的结论.

\begin{Lemma}
\label{lem:ch5-cotlar-type-inequality}
若 $T$ 是 Calderón--Zygmund 算子, 则对任意 $0<\nu<1$ 和 $f\in C_c^\infty$,
\begin{equation}
\label{eq:ch5-cotlar-type-inequality}
 T^*f(x)\le C_\nu\left(M(|Tf|^\nu)(x)^{1/\nu}+Mf(x)\right).
\end{equation}
\end{Lemma}

需要下面归功于 Kolmogorov 的结果.

\begin{Lemma}
\label{lem:ch5-kolmogorov-weak-estimate}
设 $S$ 是弱型 $(1,1)$ 算子, $0<\nu<1$, $E$ 是有限测度集. 则存在只依赖于 $\nu$ 的常数 $C$, 使得
\[
 \int_E|Sf(x)|^\nu\,dx\le C|E|^{1-\nu}\|f\|_1^\nu.
\]
\end{Lemma}

\begin{Proof}
由分布函数公式和弱型 $(1,1)$ 不等式,
\[
 \int_E|Sf|^\nu
 =\nu\int_0^\infty\lambda^{\nu-1}
 |\{x\in E:|Sf(x)|>\lambda\}|\,d\lambda
 \le\nu\int_0^\infty\lambda^{\nu-1}
 \min\left(|E|,\frac C\lambda\|f\|_1\right)d\lambda.
\]
在 $C\|f\|_1/|E|$ 处分割积分即得结论.
\end{Proof}

\hypertarget{chapter-05-page-115}{}
\begin{Proof}[引理~\ref{lem:ch5-cotlar-type-inequality} 的证明]
先在原点证明估计, 常数与截断尺度 $\varepsilon$ 无关. 令 $Q=B(0,\varepsilon/2)$, $2Q=B(0,\varepsilon)$, 并写 $f=f_1+f_2$, 其中 $f_1=f\chi_{2Q}$. 则 $Tf_2(0)=T_\varepsilon f(0)$. 若 $z\in Q$, 由标准核条件,
\[
 |Tf_2(z)-Tf_2(0)|\le CMf(0).
\]
所以
\begin{equation}
\label{eq:ch5-truncation-local-comparison}
 |T_\varepsilon f(0)|\le CMf(0)+|Tf(z)|+|Tf_1(z)|.
\end{equation}
当 $\nu=1$ 时, 对 $z\in Q$ 的三项分别用极大函数和 $T$ 的弱型 $(1,1)$ 估计测度, 得 $|T_\varepsilon f(0)|\le C(M(Tf)(0)+Mf(0))$.

当 $0<\nu<1$ 时, 对 \eqref{eq:ch5-truncation-local-comparison} 取 $\nu$ 次幂, 在 $Q$ 上积分, 除以 $|Q|$, 再取 $1/\nu$ 次幂. 引理~\ref{lem:ch5-kolmogorov-weak-estimate} 给出
\[
 \left(\frac1{|Q|}\int_Q|Tf_1(z)|^\nu\,dz\right)^{1/\nu}
 \le C|Q|^{-1}\|f_1\|_1\le CMf(0).
\]
由平移不变的同一论证并对 $\varepsilon$ 取上确界, 得 \eqref{eq:ch5-cotlar-type-inequality}.
\end{Proof}

\hypertarget{chapter-05-page-116}{}
\begin{Proof}[定理~\ref{thm:ch5-maximal-cz-bounds} 的证明]
当取 $\nu=1$ 的相应估计时, 因 $T$ 与 $M$ 都在 $L^p$, $p>1$, 上有界, $T^*$ 的强型 $(p,p)$ 有界性立即得到.

为证明弱型 $(1,1)$, 在 \eqref{eq:ch5-cotlar-type-inequality} 中取 $\nu<1$. $Mf$ 项已有所需估计. 对另一项使用二进极大算子 $M_d$ 与一般极大函数的比较. 令 $E=\{x:M_d(|Tf|^\nu)(x)>\lambda^\nu\}$, 则对 $f\in C_c^\infty$ 有 $|E|<\infty$. Calderón--Zygmund 分解的极大立方体给出
\[
 |E|\le\lambda^{-\nu}\int_E|Tf(y)|^\nu\,dy.
\]
由引理~\ref{lem:ch5-kolmogorov-weak-estimate}, $|E|\le C\lambda^{-\nu}|E|^{1-\nu}\|f\|_1^\nu$, 从而 $|E|\le C\lambda^{-1}\|f\|_1$. 这完成证明.
\end{Proof}
